\subsection{The Second Week}\label{sec:week2}

The Second Sunday belongs to John the Baptist, and in Year B the Church hears
first the prophecy he fulfilled (Isaiah 40:1--5, 9--11):

\begin{quote}
Be comforted, be comforted, my people, saith your God. Speak ye to the heart of
Jerusalem, and call to her: for her evil is come to an end, her iniquity is
forgiven: she hath received of the hand of the Lord double for all her sins. The
voice of one crying in the desert: Prepare ye the way of the Lord, make straight
in the wilderness the paths of our God. Every valley shall be exalted, and every
mountain and hill shall be made low, and the crooked shall become straight, and
the rough ways plain.\par
{\smallskip\noindent\small ---Isaiah 40:1--4\par}
\end{quote}

\noindent St Peter writes to those troubled by the Lord's delay: ``The Lord delayeth
not his promise, as some imagine, but dealeth patiently for your sake, not
willing that any should perish, but that all should return to penance.'' And St
Mark opens his Gospel with John in the desert, the voice Isaiah heard, and the
people going out to him confessing their sins. This is the week to go to
confession as they went, before the Lord comes (Section~\ref{sec:confession}). In Year A the prophet is Isaiah again, on the
rod of Jesse, and John preaches in Matthew; in Year C Baruch bids Jerusalem put
off her mourning, and Luke names the emperor, the governor and the high priests
of the year when the word of God came to John. In the Office of Readings,
Eusebius of Caesarea, the historian of the early Church, comments on the same
chapter of Isaiah.

Set beside the prophet's levelled hills a sentence of St Augustine, though he
spoke it of the centurion and not of Isaiah: ``The more humble, the more capacious, and the more
full. For the hills drive back the water, but the valleys are filled by it.''

The weekdays keep Isaiah's comfort. On Monday the desert blooms: ``And the
redeemed of the Lord shall return, and shall come into Sion with praise.'' The
Gospel tells of the paralytic let down through the tiles, to whom Jesus said,
``Man, thy sins are forgiven thee.'' On Tuesday Isaiah's shepherd gathers the
lambs with his arm, and in the Gospel the shepherd leaves the ninety-nine. On
Wednesday comes a promise to the weary:

\begin{quote}
But they that hope in the Lord shall renew their strength, they shall take wings
as eagles, they shall run and not be weary, they shall walk and not faint.\par
{\smallskip\noindent\small ---Isaiah 40:31\par}
\end{quote}

\noindent and Christ in the Gospel says, ``Come to me all you that labor and are
burdened, and I will refresh you.'' From Thursday the Gospels turn to John the
Baptist, whom the Lord calls the Elijah who was to come; on Saturday Sirach
praises Elijah, who ``stood up, as a fire, and his word burnt like a torch.''
In the Office of Readings this week the Church reads a chapter of St John of
the Cross on Monday, the
Second Vatican Council on the pilgrim Church on Tuesday, St Augustine on
Wednesday, St Peter Chrysologus on Thursday, St Irenaeus on Friday and Isaac of
Stella, a Cistercian abbot, on Saturday.

\subsubsection*{The Immaculate Conception}

On 8 December the Church keeps the solemnity of the Immaculate Conception of
the Blessed Virgin Mary, patroness of the United States. It is a holy day of
obligation in this country by the decree of the bishops confirmed by the Holy
See in 1992 and in force since 1 January 1993; the rule that lifts the
obligation when a holy day falls on a Saturday or a Monday reaches only 1
January, 15 August and 1 November, never 8 December. When 8 December falls on a
Sunday, as it will next in 2030, the solemnity is kept on the Monday, and the
obligation to attend Mass does not pass to the Monday with it (Dicastery for
Divine Worship and the Discipline of the Sacraments, Note of 23 January 2025,
n.~7).

The first reading is the curse of the serpent, read with Genesis 3:20, in which
Adam names his wife Eve, the mother of the living:

\begin{quote}
I will put enmities between thee and the woman, and thy seed and her seed: she
shall crush thy head, and thou shalt lie in wait for her heel.\par
{\smallskip\noindent\small ---Genesis 3:15\par}
\end{quote}

\noindent The Gospel is the Annunciation, with the angel's greeting, ``Hail, full
of grace, the Lord is with thee: blessed art thou among women.'' When Blessed
Pius IX defined the doctrine on 8 December 1854, in the apostolic constitution
\emph{Ineffabilis Deus}, he gathered how the Fathers had read these two texts. Of the first, he
said that ``they taught that in this Divine oracle was clearly and openly
pointed out the merciful Redeemer of the human race,'' and his Mother with him.
Then came the definition:

\begin{quote}
We declare, pronounce, and define that the doctrine which holds that the Blessed
Virgin Mary, at the first instant of her conception, by a singular privilege and
grace of the Omnipotent God, in virtue of the merits of Jesus Christ, the
Saviour of mankind, was preserved immaculate from all stain of original sin, has
been revealed by God, and therefore should firmly and constantly be believed by
all the faithful.\par
{\smallskip\noindent\small ---Blessed Pius IX, \emph{Ineffabilis Deus}\par}
\end{quote}

The Church in this country had chosen her patroness eight years before. In May
1846, at the Sixth Provincial Council of Baltimore, the bishops acted
together: \latin{Patres, ardentibus votis, plausu consensuque unanimi, B. Mariam
Virginem sine labe Conceptam in Patronam Americæ Fœderatæ Statuum elegerunt},
``The Fathers, with ardent desire, with applause and unanimous consent, chose
the Blessed Virgin Mary conceived without stain as Patroness of the United States
of America.'' Pius IX confirmed the choice the following year. In the Office
of Readings on 8 December the Church reads a prayer of St Anselm to our Lady,
and ends with the \latin{Te Deum}.

Through the week of the Immaculate Conception a household can end the day as
the Roman Breviary ended Compline, its Night Prayer, all through Advent: with
the antiphon \latin{Alma Redemptoris Mater}, ``Loving Mother of the Redeemer,''
which greets her as the open gate of heaven and the star of the sea, and asks
her who bore her own holy Maker, while nature marvelled, to have pity on
sinners. It is given in Latin and English in Appendix~\ref{sec:prayers}.

\subsubsection*{The saints of the week}

On 9 December the Church keeps St Juan Diego, the Mexican to whom our Lady
showed herself at Guadalupe, and on 10 December, since 2019, Our Lady of
Loreto. The decree that added Loreto to the calendar gave its reasons: the Holy
House is a memorial of the mystery of the Incarnation, it calls to mind the
virtues of the Holy Family, it gave its name to the Litany of Loreto, and Our
Lady of Loreto is the patroness of those who fly. On 11 December comes Pope St
Damasus, of whom the Roman Martyrology of 1902 said,
\latin{multa etiam sanctorum Martyrum corpora invenit, eorumque memorias versibus
exornavit}, ``he also found the bodies of many holy martyrs, and adorned their
memorials with verses.''

On the Friday of this week the Office of Readings reads St Irenaeus on the Virgin who undid the disobedience of the first virgin:

\begin{quote}
And if the former did disobey God, yet the latter was persuaded to be obedient
to God, in order that the Virgin Mary might become the patroness (advocata) of
the virgin Eve. And thus, as the human race fell into bondage to death by means
of a virgin, so is it rescued by a virgin; virginal disobedience having been
balanced in the opposite scale by virginal obedience.\par
{\smallskip\noindent\small ---St Irenaeus, \emph{Against Heresies}, V.19\par}
\end{quote}

\noindent Elsewhere in the same work, in a passage the Office does not read in
Advent, Irenaeus gathered it into one sentence: ``For what the virgin Eve had bound
fast through unbelief, this did the virgin Mary set free through faith.''

On 12 December the Church in the United States keeps the feast of Our Lady of
Guadalupe, with readings from Zechariah or from Revelation on the woman
clothed with the sun, and the Gospel of the Annunciation or the Visitation. The
\emph{Directory on Popular Piety and the Liturgy} observes that this feast
stands at the centre of the popular piety of the Americas in the days before
Christmas (n.~102). The Nahuatl account of the apparitions, the \emph{Huei
tlamahuiçoltica}, gives in Juan Diego's own tongue what our Lady said to him when
he feared for his dying uncle: \textit{cuix àmo nican nicà nimoNantzin?}, ``Am I
not here, I who am thy mother?'' \textit{cuix àmo noçehuallotitlan,
nècauhyotitlan in ticà?}, ``Art thou not in my shadow and under my shelter?''
